\documentclass[12pt]{article}
\usepackage[a4paper, margin=2cm]{geometry}

\usepackage{fontspec}
\usepackage{polyglossia}
\setdefaultlanguage{russian}
\setmainlanguage{russian}
\setmainfont{Times New Roman}

\begin{document}

\begin{center}
\textbf{\LARGE СЛУЖЕБНАЯ ЗАПИСКА}
\end{center}

\vspace{1cm}

Кому: Иван Иванов, начальник отдела по международному экспорту нефти

От: Александр Александрович, сотрудник отдела экспорта, Санкт-Петербург

Суть вопроса: предлагаю скорректировать подход к планированию поставок.

Уважаемый Иван Иванов,

В настоящее время расчеты плановых показателей ведутся на основании усреднённых цен, что приводит к расхождениям между плановыми и фактическими поставками и повышает риски финансового характера. Предлагаю ориентироваться на актуальную рыночную стоимость нефти Brent, чтобы повысить точность расчетов и снизить риски.

Обоснование: Brent демонстрирует значительную волатильность; по данным на 25 марта 2026 года текущая цена Brent составляет около 102 USD за баррель. Источники: Trading Economics, OilPrice.com, Business Insider, CNBC, FT Markets.

Предлагаю рассмотреть следующие шаги по внедрению данного подхода:
- выбрать единый источник цен для расчета;
- наладить ежедневное обновление цен;
- запустить пилотный проект в одном из направлений поставок на ближайший квартал;
- оценить влияние на точность планирования и финансовые риски.

С уважением,
Александр Александрович
Сотрудник отдела экспорта, Санкт-Петербург

\vfill

\noindent
Дата: 25.03.2026 \hfill Подпись: \rule{5cm}{0.4pt}

\end{document}